\documentclass{article}
\usepackage[T1]{fontenc}
\usepackage[utf8]{inputenc}
\usepackage{lmodern}
\usepackage[margin=0.8in]{geometry}
\usepackage{amsmath,amssymb}
\usepackage[dvipsnames]{xcolor}
\usepackage{booktabs}
\usepackage{fancyvrb}
\usepackage{parskip}
\usepackage[colorlinks=true,linkcolor=Maroon,urlcolor=Blue,pdftitle={CS3423 Mini Assignment 6: Q5}]{hyperref}

\setcounter{secnumdepth}{-\maxdimen} % remove section numbering
\providecommand{\tightlist}{\setlength{\itemsep}{0pt}\setlength{\parskip}{0pt}}

\title{CS3423 Mini Assignment 6: Q5}
\author{}
\date{}

\begin{document}
\maketitle

\section{Q5: mem2reg Pass Analysis}

\subsection{Setup}

I generated \texttt{before.ll} and \texttt{after.ll} for each of the 5
programs using the commands (given in README):

\begin{Verbatim}[fontsize=\small]
clang++ -std=c++17 -S -emit-llvm -O0 -Xclang -disable-O0-optnone "cpp Codes/N.cpp" -o IR/N_before.ll
opt -passes=mem2reg -S IR/N_before.ll -o IR/N_after.ll
opt -passes=verify -disable-output IR/N_after.ll
\end{Verbatim}

The flag \texttt{-disable-O0-optnone} is needed because clang normally
adds the \texttt{optnone} attribute to functions compiled with
\texttt{-O0}. This attribute tells LLVM to skip most optimization passes
for that function, including mem2reg.

\subsection{(a) What analysis information does mem2reg depend on?}

The mem2reg pass uses information about the control flow graph and the
definitions and uses of variables to convert eligible variables stored
in memory into SSA form.

\medskip
\noindent
\begin{tabular}{@{}p{0.30\linewidth}p{0.66\linewidth}@{}}
\toprule
Analysis & What it provides / how it is used \\
\midrule
DominatorTreeAnalysis & Provides the dominator tree, used to compute where
PHI nodes are needed and to rename values during SSA construction. \\
AssumptionAnalysis & Provides \texttt{llvm.assume} information, used so
that assumes on a promoted value are handled correctly. \\
TargetIRAnalysis & Provides target-specific information about the IR. It
is run by the pass manager for \texttt{PromotePass} (it appears in the
\texttt{-debug-pass-manager} trace). \\
\bottomrule
\end{tabular}
\medskip

\begin{itemize}
\tightlist
\item \textbf{Iterated dominance frontier:} mem2reg computes this from the
  dominator tree to decide where PHI nodes must be placed for each
  promoted variable. It is not a separate analysis pass.
\item \textbf{Promotability check (definitions and uses):} an alloca is
  promoted only if it has a single-value type (not an array or struct)
  and every use is a non-volatile load or store that uses the alloca as
  its pointer operand. If its address is passed to a call, stored
  somewhere, used in a GEP, or fed into a PHI/select, it is not promoted.
\end{itemize}

Using this information, mem2reg inserts PHI nodes where required and
renames variables so that each SSA value has a single definition.

\subsection{(b) What variables are converted to SSA form?}

\medskip
\noindent
\begin{tabular}{@{}c l c c c c@{}}
\toprule
Program & Function & Allocas before & Allocas after & PHIs before & PHIs after \\
\midrule
1 & basicPointer & 3 & 0 & 0 & 0 \\
2 & conditionalPointer & 4 & 2 & 0 & 1 \\
3 & pointerEscape + modify & 4 & 1 & 0 & 0 \\
4 & arrayPointer & 5 & 1 & 1 & 3 \\
5 & heapPointer & 4 & 0 & 0 & 2 \\
\bottomrule
\end{tabular}
\medskip

\subsubsection{Program 1: basicPointer}

The variables \texttt{input}, \texttt{value} and \texttt{ptr} are
promoted to SSA form. Their corresponding alloca, load and store
instructions are eliminated.

Even though the address of \texttt{value} is initially stored in
\texttt{ptr}, both variables can be promoted. Once \texttt{ptr} is
promoted, its uses become direct uses of the address of \texttt{value}.
This allows \texttt{value} to be promoted as well. The resulting IR uses
SSA values directly instead of accessing these variables through memory.

\subsubsection{Program 2: conditionalPointer}

The variables \texttt{condition} and \texttt{ptr} are promoted to SSA
form. Their corresponding stack allocations and associated loads and
stores are removed.

A pointer-valued PHI node is introduced at the join block. It selects
between the address of \texttt{first} and the address of
\texttt{second}, depending on which branch was executed. The selected
pointer is then used to access and update the corresponding integer.

\subsubsection{Program 3: pointerEscape}

In \texttt{pointerEscape()}, the variables \texttt{n} and \texttt{ptr}
are promoted to SSA form. In \texttt{modify()}, the stack allocation
used to store the parameter \texttt{p} is also eliminated because the
parameter is already available as an SSA value.

The corresponding alloca, load and store instructions for these promoted
variables are removed. However, the actual variable \texttt{value}
remains memory-backed because its address is passed to
\texttt{modify()}.

\subsubsection{Program 4: arrayPointer}

The scalar variables \texttt{n}, \texttt{ptr}, \texttt{sum} and
\texttt{i} are promoted to SSA form. Their corresponding stack
allocations and associated loads and stores are removed.

PHI nodes are introduced for the loop-carried variables \texttt{sum} and
\texttt{i}, since their values are updated during each loop iteration.
The IR also contains a PHI node associated with the short-circuit
evaluation of the \texttt{\&\&} condition.

\subsubsection{Program 5: heapPointer}

The local stack variables \texttt{n}, \texttt{buffer}, \texttt{sum} and
\texttt{i} are promoted to SSA form. Their corresponding stack
allocations and associated loads and stores are eliminated.

Two PHI nodes are introduced at the loop header for \texttt{sum} and
\texttt{i}, since their values are carried from one iteration to the
next. The pointer variable \texttt{buffer} is also promoted, but the
dynamically allocated array it points to remains in heap memory.

\subsection{(c) What variables are NOT converted to SSA form and why?}

\subsubsection{Program 1}

No variables remain unpromoted. All three local variables
(\texttt{input}, \texttt{value} and \texttt{ptr}) are successfully
converted to SSA form.

\subsubsection{Program 2}

The variables \texttt{first} and \texttt{second} remain as allocas. Their
addresses are used as incoming values of the pointer PHI node, which is a
use other than a direct load or store. They therefore fail the
promotability check, and the loads and stores through the selected
pointer stay in memory.

\subsubsection{Program 3}

The variable \texttt{value} remains as an alloca in memory. Its address
is passed to the function \texttt{modify()}, which accesses and modifies
the value through the pointer. Therefore, mem2reg cannot promote this
variable in this example.

\subsubsection{Program 4}

The array \texttt{arr} remains memory-backed. It has an array type rather
than a single-value type, and its elements are accessed through GEP
instructions, which are not plain loads or stores of the alloca.
mem2reg does not promote such allocas (splitting an aggregate into
scalars is the job of SROA, not mem2reg).

\subsubsection{Program 5}

The dynamically allocated array created using \texttt{new int[n]}
remains in heap memory. mem2reg operates on eligible stack allocations
created using \texttt{alloca}, not dynamically allocated heap objects.

The local pointer variable \texttt{buffer}, which stores the address of
this heap allocation, is promoted to SSA form. However, the actual heap
array and its load/store operations remain unchanged.

\subsection{Conclusion}

The mem2reg pass converts promotable stack variables into SSA values,
removing their alloca, load and store instructions. It uses the dominator
tree to place PHI nodes where values from different control-flow paths
merge. Allocas whose address escapes (passed to a call, stored, or used
in a PHI or GEP), aggregate allocas such as arrays, and heap objects are
not promoted and remain memory-backed.

\end{document}